\documentclass[10pt,a4paper]{amsart}
\usepackage[margin=19mm]{geometry}
\usepackage{amsmath,amssymb,mathtools}
\usepackage[scr=rsfs,cal=euler]{mathalfa}
\usepackage{xfrac}
\usepackage[unicode,hidelinks,bookmarks=false]{hyperref}
\newcommand{\ie}{i.e.\ }
\newcommand{\cf}{cf.\ }
\newcommand{\resp}{respectively\ }
\newcommand{\Subsection}{\subsection}
\providecommand{\nobreakdash}{\nobreak\hbox{-}}
\setlength{\emergencystretch}{2em}
\multlinegap0pt
\usepackage{bm}
\usepackage{bbm}
\usepackage{dsfont}
\usepackage{xspace}
\usepackage{cancel}
\usepackage{mathtools}
\usepackage{amsthm}
\makeatletter
\@ifundefined{theorem}{\newtheorem{theorem}{Theorem}}{}
\makeatother
\def\Gp{{G_p}}
\title[On the restriction of some irreducible mod $p$ representatio]{On the restriction of some irreducible mod $p$ representations}
\author{Shubhanshi Gupta}
\date{Source: arXiv:2405.20007v3 (CC BY 4.0)}
\begin{document}
\maketitle

\noindent\emph{Source and licence.} On the restriction of some irreducible mod $p$ representations. Version arXiv:2405.20007v3 (CC BY 4.0), distributed under the Creative Commons Attribution 4.0 International licence, as recorded in the arXiv licence metadata for this exact version and shown on the abstract page https://arxiv.org/abs/2405.20007v3. Full rights notice: licenses/arxiv-2405.20007-CC-BY-4.0.txt.\par\smallskip

\noindent\textit{Editorial scope.} Counted unit: the Theorem whose source label is \texttt{T3} in the article (source0.tex lines 453--470, source label \texttt{T3}), delivered together with the article's own complete original proof of it (source0.tex lines 472--503, one \texttt{proof} environment of 32 lines). No mathematics is written, repaired or re-derived: statement and proof are the article's own text line for line. Required context retained uncounted: T2 (lines 320--322). Every definition, display and label of those lines is retained unchanged. Omitted from this package and not redistributed: 1--319, 323--452, 471--471, 504--561. Nothing inside a retained excerpt is omitted or altered: every retained body is delivered exactly as the article's lines are written, so no omission receipt of any kind accompanies this package. Citations: the retained excerpts carry no citation command at all, so no bibliography entry of the article is transcribed and no reference is redistributed. Reference adaptation: none was needed and none was made; every reference inside the delivered unit resolves inside it, so no unresolved reference or citation is produced. Printed slips kept literal: no printed word, symbol, hypothesis or number of the counted statement or of its proof was altered. Layout and class: the article's own class is \texttt{amsart} with packages including amscd, amssymb, latexsym, amsmath, comment, hyperref, xy, anysize, mathpazo, bbm, color, dsfont, mathtools, cases; the wrapper is the lane's pinned \texttt{amsart} editorial wrapper at 10pt on A4, which restates the theorem-like environments the retained text needs and adds the article's own macro definitions that the retained text needs (\texttt{\textbackslash Gp}), with extra wrapper packages (bm, bbm, dsfont, xspace, cancel, mathtools). The delivered numbering is the unit's own. Assets: no retained span contains a figure environment, a graphics-inclusion command, a PDF-page inclusion, a table or a TikZ picture; no image file is copied into this package, the package declares no asset dependency and no image file is redistributed with it. Licence: version arXiv:2405.20007v3 of this article is distributed under the Creative Commons Attribution 4.0 International licence, as recorded in the arXiv licence metadata for this exact version and shown on the abstract page https://arxiv.org/abs/2405.20007v3. A full rights notice is delivered at licenses/arxiv-2405.20007-CC-BY-4.0.txt. Characters: every retained excerpt is pure ASCII, so the delivered engine, XeLaTeX with its default Latin Modern text font, reports no missing character.
	\begin{theorem}\label{T2}
		Let $\vec{r}=(r_1,r_2)$ with $0\leq r_2\leq r_1< p-r_2.$ Then $${V_{\vec{r}}|}_{\Gp}\cong V_{r_1+r_2}\oplus V_{r_1+r_2-2}(1)\oplus \cdots\oplus V_{r_1-r_2}(r_2).$$
	\end{theorem}

	\begin{theorem}\label{T3}
		Let $\vec{r}=(r_1,r_2,r_3)$ such that $ 0\leq r_3\leq r_2\leq r_1< p-r_2-r_3.$ Then
		\begin{enumerate}
			\item if $r_1\geq r_2+r_3,$
			\begin{equation*}
			    {V_{\vec{r}}|}_{\Gp}\cong            	\displaystyle{\left[\bigoplus\limits_{i=0}^{r_3}(i+1)\oplus\bigoplus\limits_{i=r_3+1}^{r_2}(r_3+1)\oplus \bigoplus\limits_{i=r_2+1}^{r_2+r_3}(r_3+(r_2+1)-i)\right]V_{r_1+r_2+r_3-2i}(i),}
			\end{equation*}				
			
			\item if $r_1< r_2+r_3,$
			\begin{multline*}
			    {V_{\vec{r}}|}_{\Gp}\cong 
				\left[\displaystyle{\bigoplus\limits_{i=0}^{r_3}(i+1)\oplus\bigoplus\limits_{i=r_3+1}^{r_2}(r_3+1)\oplus\bigoplus\limits_{i=r_2+1}^{r_1}(r_3+(r_2+1)-i)\oplus}\right.\\
\left.\displaystyle{\bigoplus\limits_{i=r_1+1}^{\left\lfloor\frac{r_1+r_2+r_3}{2}\right\rfloor}(r_1+r_2+r_3-2i+1)}\right] V_{r_1+r_2+r_3-2i}(i),
			\end{multline*}
where $\lfloor \cdot\rfloor$ denotes the floor function.
			
		\end{enumerate}
	\end{theorem}

	\begin{proof}
		We will prove this theorem by induction. Note that for $r_3=0$, $V_{r_1}\otimes V_{r_2}\otimes V_0\cong V_{r_1}\otimes V_{r_2}$ and so by Theorem \ref{T2}, the result is true in this case. With the induction hypothesis that the theorem holds for $r_3-1\leq r_2-1\leq r_1$, we will show that it holds for $r_3\leq r_2\leq r_1$. (So to decompose $V_{r_1}\otimes V_{r_2}\otimes V_{r_3},$ we can apply induction $r_3$ times on the initial case $V_{r_1}\otimes V_{r_2-r_3}\otimes V_{0}.$) Note that,
		\begin{align*}
			V_{r_1}\otimes V_{r_2}\otimes V_{r_3}&\cong V_{r_1}\otimes \left(\displaystyle{\bigoplus\limits_{i=0}^{r_3}V_{r_2+r_3-2i}(i)}\right)~~~~~(\text{By Theorem }\ref{T2})\\
			&\cong (V_{r_1}\otimes V_{r_2+r_3})\oplus (V_{r_1}\otimes (V_{r_2-1}\otimes V_{r_3-1})\otimes \mathrm{det}).
		\end{align*}\\
        
		Using induction hypothesis and Theorem \ref{T2},
		\begin{enumerate}
			\item when $r_1\geq r_2+r_3,$            
			\begin{multline*}			    
			V_{r_1}\otimes V_{r_2}\otimes V_{r_3}\cong \displaystyle{\bigoplus\limits_{i=0}^{r_2+r_3}V_{r_1+r_2+r_3-2i}(i)\oplus}\left[\left(\displaystyle{\bigoplus\limits_{i=0}^{r_3-1}(i+1)\oplus \bigoplus\limits_{i=r_3}^{r_2-1}(r_3)\oplus}\right.\right.\\
\left.\left.\displaystyle{\bigoplus\limits_{i=r_2}^{r_2+r_3-2}(r_3-1+(r_2)-i)}\right)V_{r_1+r_2-1+r_3-1-2i}(i)\otimes \mathrm{det}\right]		
            \end{multline*}
			$\displaystyle{\cong\bigoplus\limits_{i=0}^{r_2+r_3}V_{r_1+r_2+r_3-2i}(i)\oplus\left[\left(\bigoplus\limits_{i=1}^{r_3}i\oplus \bigoplus\limits_{i=r_3+1}^{r_2}r_3\oplus \bigoplus\limits_{i=r_2+1}^{r_2+r_3-1}(r_3+r_2-i)\right)V_{r_1+r_2+r_3-2i}(i)\right]}$\\			
$\displaystyle{\cong \left(\bigoplus\limits_{i=0}^{r_3}(i+1)\right.}\oplus\displaystyle{\left.\bigoplus\limits_{i=r_3+1}^{r_2}(r_3+1)\oplus\bigoplus\limits_{i=r_2+1}^{r_2+r_3}(r_3+(r_2+1)-i)\right)}V_{r_1+r_2+r_3-2i}(i),$\\
              
			
			\item when $r_1<r_2+r_3,$
			\begin{multline*}
			    V_{r_1}\otimes V_{r_2}\otimes V_{r_3}\cong \displaystyle{\bigoplus\limits_{i=0}^{r_1}V_{r_1+r_2+r_3-2i}(i)\oplus}\left[\mathrm{det}\otimes\left(\displaystyle{\bigoplus\limits_{i=0}^{r_3-1}(i+1)\oplus \bigoplus\limits_{i=r_3}^{r_2-1}r_3\oplus\bigoplus\limits_{i=r_2}^{r_1}(r_3-1+(r_2)-i)\oplus}\right.\right.\\		
\displaystyle{\left.\left.\bigoplus\limits_{i=r_1+1}^{\left\lfloor\frac{r_1+r_2+r_3-2}{2}\right\rfloor}(r_1+r_2+r_3-2i-1)\right)V_{r_1+r_2+r_3-2i-2}(i)\right]}			\end{multline*}
            \begin{multline*}
			    \displaystyle{\cong \bigoplus\limits_{i=0}^{r_1}V_{r_1+r_2+r_3-2i}(i)\oplus \left[\left(\bigoplus\limits_{i=1}^{r_3}i\oplus \bigoplus\limits_{i=r_3+1}^{r_2}r_3\oplus\bigoplus\limits_{i=r_2+1}^{r_1+1}(r_3+r_2-i)\oplus\right.\right.}\\
\displaystyle{\left.\left.\bigoplus\limits_{i=r_1+2}^{\left\lfloor\frac{r_1+r_2+r_3}{2}\right\rfloor}(r_1+r_2+r_3-2i+1) \right)V_{r_1+r_2+r_3-2i}(i)\right]}            
			\end{multline*}
            \begin{equation*}
        \displaystyle{\cong\left(\bigoplus\limits_{i=0}^{r_3}(i+1)\oplus\bigoplus\limits_{i=r_3+1}^{r_2}(r_3+1)\oplus\bigoplus\limits_{i=r_2+1}^{r_1}(r_3+(r_2+1)-i)\oplus\bigoplus\limits_{i=r_1+1}^{\left\lfloor\frac{r_1+r_2+r_3}{2}\right\rfloor}(r_1+r_2+r_3-2i+1)\right)}V_{r_1+r_2+r_3-2i}(i).
            \end{equation*}
		\end{enumerate}
		This proves the theorem.
	\end{proof}

\end{document}
